\documentclass[11pt]{article}
\usepackage[margin=1in]{geometry}
\usepackage{amsmath}
\usepackage{booktabs}

\title{Random Forests Solve Two Moons: Non-Linear Boundaries Without Feature Engineering}
\author{Paper2Proof canary set\thanks{A small paper written for the Paper2Proof evaluation set (eval set A).}}
\date{2026}

\begin{document}
\maketitle

\begin{abstract}
The two-moons problem is the textbook example of a non-linear decision boundary. A random forest
of 200 trees outperforms logistic regression on noisy two moons by 0.061 accuracy and reaches a
test accuracy of 0.990.
\end{abstract}

\section{Introduction}
Linear models cannot separate interleaved half circles. We quantify the gain of a random forest
over logistic regression on a fixed, noisy two-moons sample.

\section{Method}
\paragraph{Random forest.} 200 trees, Gini impurity, at least 5 samples per leaf, bootstrap
sampling, $\sqrt{n_\text{features}}$ candidate features per split, random state equal to the seed.
\paragraph{Baseline.} Logistic regression with an L2 penalty ($C = 1.0$) and the L-BFGS solver.

\section{Experimental setup}
\paragraph{Data.} One fixed sample of 1000 points from scikit-learn's \texttt{make\_moons} with
Gaussian noise of standard deviation 0.3, generated with \texttt{random\_state=0}.
\paragraph{Splits.} For each seed $s \in \{0, 1, 2\}$, a stratified split of that sample with
30\,\% test points (\texttt{train\_test\_split(test\_size=0.3, random\_state=s, stratify=y)});
both models use the same split.
\paragraph{Metric.} Test accuracy, mean and sample standard deviation over the three seeds; the
gain is the per-seed difference averaged over seeds.
\paragraph{Software.} Python 3.12, scikit-learn 1.9.1, CPU only.

\section{Results}
Table~\ref{tab:moons} reports both models. The random forest outperforms logistic regression by
0.061 accuracy, and it reaches a test accuracy of 0.990 $\pm$ 0.004.

\begin{table}[t]
\centering
\caption{Test accuracy on two moons, noise 0.3 (mean $\pm$ standard deviation over seeds 0, 1, 2).}
\label{tab:moons}
\begin{tabular}{lc}
\toprule
Model & Accuracy \\
\midrule
Random forest (200 trees, min. 5 samples per leaf) & $0.990 \pm 0.004$ \\
\midrule
Gain over logistic regression & $+0.061$ \\
\bottomrule
\end{tabular}
\end{table}

\section{Conclusion}
A random forest captures the two-moons boundary with no feature engineering.

\end{document}
